\subsection*{如何使用本书}
\addcontentsline{toc}{subsection}{如何使用本书}
建议先按章节阅读主文，再回看对应 lecture workspace 里的 \texttt{lecture\_plan.json}、\texttt{coverage\_units.jsonl}、\texttt{figure\_manifest.json} 和 \texttt{eval\_reports/pass\_\#\#.json}。如果你只是把这本书当教材来读，可以直接沿章节顺序通读；如果你把它当研究资料来用，则应同时查看每章的 source 与 coverage 侧车。

如果某章尚未进入 \texttt{book/}，优先看该章的 omission log 与 repair log，再决定是补源、修复，还是接受缺口。对于已经进入本书的章节，也仍建议把 omission log 当作正文的附属阅读，因为它会告诉你哪些地方是“证据不足但已显式处理”的，而不是“完全不存在问题”。

为避免逐讲原始版面中的局部页码与合并后的全书页码混淆，all-in-one PDF 会在每一页右上角额外标出统一连续页码；全书引用与翻页时请以该统一页码为准。

若要引用本书中的具体论断、图或代码说明，推荐采用三层引用方式：
\begin{itemize}
    \item 第一层：引用本书的章节标题与统一页码；
    \item 第二层：引用该章对应的官方课程页或官方视频；
    \item 第三层：必要时回到工作区中的 \texttt{figure\_manifest.json}、\texttt{source\_manifest.json} 与 references 条目，给出更精确的原始出处。
\end{itemize}

简言之：本书适合连续学习，也适合带着审计意识来读。正文负责教学组织，侧车负责来源与边界条件。
